\section{Evaluation}
\label{sec:eval}

The evaluation asks how far structure transfer reaches the earlier cohort in realistic settings
(\cref{sec:eval-transfer,sec:eval-real,sec:eval-induction}), whether the conditions of
\cref{thm:main,thm:reactive} hold in the compiled engines and which are load-bearing
(\cref{sec:eval-game,sec:eval-reactive,sec:eval-audit}), which channels still reach cohort H
(\cref{sec:eval-cohorth,sec:eval-rw}), and what the designs cost
(\cref{sec:eval-perf,sec:eval-recall}).

\paragraph{Setup}
Records are embedded with \texttt{bge-base-en-v1.5} \cite{xiao2024cpack}. Corpora are D2, eight
BEIR domains \cite{thakur2021beir} used as a topic taxonomy; D3, a synthetic procedural-memory
corpus; 13{,}894 timestamped \texttt{kubernetes} issues labeled by their \texttt{sig/*} area; 20
Newsgroups; and memories a local model wrote from those issues. Advantage is accuracy minus chance,
untruncated, for an adversary that guesses wherever it has no label. Intervals are 95\%, over each
experiment's independent units: bootstrap from twenty units, Student-$t$ below; the supplementary
material names the unit and its count for every experiment and specifies the learned attacker.
Attacks read only the engine's emitted event stream, which is the storage adversary's view, and
every number in the text is generated from the stored results. Outcomes registered before a run are
marked as such, including three that failed (\cref{sec:eval-induction}, the composed attack of
\cref{sec:eval-multisnapshot}, and query equality at $s = 4$ in the supplementary material).

\subsection{Structure Transfer, and What Drift Is Worth}
\label{sec:eval-transfer}
\label{sec:eval-drift}

The estimators are frequency analysis and an expectation-maximization estimator, neither of which
sees the true partition of an induced record; induced records are excluded when scoring, and the
intersection oracle is retained only as an upper bound.

\begin{figure}[t]
\centering
\includegraphics[width=\linewidth]{figures/fig_drift.pdf}
\caption{What subject drift is actually worth. The horizontal axis is the time gap between cohort H
and the content the adversary induces with; bands are 95\% Student-$t$ intervals over five
$k$-means seeds. An adversary that supplies archival content from cohort H's own era recovers the
advantage drift appears to remove.}
\label{fig:drift}
\end{figure}

On the timestamped \texttt{kubernetes} corpus, a chronological split lowers the advantage on the
earlier cohort substantially, which is the protection drift is supposed to give. \Cref{fig:drift}
shows what that leaves open: for a public tracker the adversary can induce with archival content,
and the advantage depends on the time gap between cohort H and the induced content. Drift protects
early records against an adversary that induces with current material, and considerably less
against one that does not. The coarse-taxonomy arm has chance one in four rather than one in twelve.

Two features of the curve are systematic. The advantage turns \emph{negative} at the largest gap
(\REfourNaive, an interval below zero): the attacker commits to a partition-to-topic map fitted on
the induced era, and once cluster geometry has drifted that map mislabels cohort H more often than
a uniform guess. The gap-zero advantage (\REfourAware) exceeds the i.i.d.\ shuffled control
(\REfourShuffled) because the shuffle mixes every era and dilutes the era-specific partition
structure that same-era induction matches; the comparison that isolates drift is gap zero against
the far eras. The supplementary material also reports advantage over the majority-class baseline.

\subsection{Deployed Stores, and Time-Placed Controls}
\label{sec:eval-real}

The attack was run against four deployed stores (Chroma, FAISS, LanceDB and sqlite-vec), an indexed
LanceDB variant and an append-only control, each inside a \texttt{gocryptfs} volume, whose
ciphertext mirrors the plaintext file structure block for block. Block labels are fitted on one
half of the induced writes and scored on the other, and induced writes are interleaved across
topics so that tail-appending stores cannot make write order masquerade as content placement.

Transfer was not detected in any of the six configurations at this budget: held-out localisation
sits at chance, and no interval on \advhist{} reaches above $\REsixWorstHi$ over \REsixSeeds{}
seeds. In the time-placed stores, and in those that rewrite their whole index on every save, an
induced write touches none of the blocks that hold cohort-H records. Chroma's HNSW insert does touch
cohort-H blocks, yet the labels carry no topic, \advhist{} $= \REsixChroma$, because which blocks an
insert touches depends on graph structure and insertion order. These stores establish the premise of
\cref{lem:transfer}, not its conclusion: transfer needs a layout that discloses which partition
holds a record at the observer's granularity.

\paragraph{A deployed store the attack reaches: pgvector}
PostgreSQL's \texttt{pgvector} with an \texttt{ivfflat} index, built on the cohort and inserted into
afterward through real SQL with a checkpoint after each write, in the same \texttt{gocryptfs}
volume, assigns an insert to its nearest list and writes that list's pages. The attack reaches
cohort H: held-out \advhist{} is \REsixteenPgAdv{} over \REsixteenPgSeeds{} seeds against chance
\REsixteenPgChance, the interval above zero (registered before the run). It is weaker than the
growing layout, because write-ahead-log and heap churn blur list pages at 4\,KiB, but it is a
deployed, content-placed store whose placement transfers to the pre-compromise cohort. Structure
transfer is therefore neither universal nor hypothetical: most deployed stores tested do not
expose a content-disclosing layout at block granularity, and one widely used store does.

\subsection{Real-Agent Induction and the Deduplication Oracle}
\label{sec:eval-induction}
\label{sec:eval-dedup}

\begin{table}[t]
\centering
\caption{At the largest budget measured (32 turns per topic); the supplement gives every budget. The attack on 3,666 memories a local model actually wrote (10 topics, chance 0.100). $w_{\mathrm{real}}$ is the burst size: how many same-topic memories one writing turn contributes to a flush. The growing layout discloses each record's partition in the snapshot and is the positive control; the advantage is scored for an adversary that guesses where it has no label.}
\label{tab:e2}
\small
\setlength{\tabcolsep}{4pt}
\begin{tabular}{lrrr}
\toprule
Layout & $w_{\mathrm{real}}$ & Turns/topic & \advhist{} [95\% CI] \\
\midrule
growing (control) & 1 & 32 & 0.051\,[0.017, 0.091] \\
growing (control) & 2 & 32 & 0.158\,[0.119, 0.191] \\
growing (control) & 4 & 32 & 0.188\,[0.173, 0.203] \\
Mnemos-P & 1 & 32 & 0.009\,[0.000, 0.018] \\
Mnemos-P & 2 & 32 & 0.014\,[0.009, 0.020] \\
Mnemos-P & 4 & 32 & 0.021\,[0.014, 0.028] \\
\bottomrule
\end{tabular}
\end{table}


\Cref{tab:e2} runs the attack against memories a local model authored from adversary-supplied
issues, and sweeps the burst size: a writing turn yields several memories on one subject within
milliseconds, so a realistic flush carries a small same-topic burst. The growing layout is the
positive control: labels learned from these memories transfer to the earlier cohort at every budget,
so agent-written text keeps the structure the attack needs. Against Mnemos-P the harness supplies
cohort-H membership, as a version-diff adversary would obtain it, isolating what padding does to
labelling. A single-record attack stays near chance, no interval above $+\REtwoSingleMaxHi$, but a
burst multiplies the signal against the fixed cover: at four records per turn the mean advantage is
\REtwoBurstRatio{} times the single-record value, peaking at $+\REtwoBurstMax$. We registered that
the padded design would stay within 0.05 of chance throughout; it does not.

The attack is cheap in the adversary's currency. An induced turn is one supplied document or
ticket, and the agent spends \REtwentyeightSecLo{} to \REtwentyeightSecHi~s on it with these models.
On the growing layout the advantage clears zero at the smallest budget tested,
\REtwentyeightGrowOneMsgs{} messages over ten topics (about \REtwentyeightGrowOneMin~min of agent
time); on Mnemos-P a single-record attack needs \REtwentyeightPOneMsgs{} messages, and a burst of two
clears zero at \REtwentyeightPTwoMsgs, though at a small advantage (descriptive analysis E28 of
the E2 logs).

Each interval resamples \emph{independent induction sequences}, each drawing its own cohort split,
induction order and attacker fit. The longitudinal study (supplementary material) reports different
advantages at nominally similar budgets because its budget counts induced records per topic while
this one counts turns, each a burst of $w_{\mathrm{real}}$ records; in records per topic the two
agree.

The deduplication oracle of \cref{sec:attack-derived} observes only \emph{whether a write happened}
after the adversary supplies a document. With a similarity threshold deciding, the oracle is
perfect. With a model deciding, as in mem0's pipeline, it is noisy but open:
\texttt{\REthreeModel} writes \REthreeLlmWriteMember{} of facts the store already holds against
\REthreeLlmWriteNonmember{} of new ones, an AUC of \REthreePLlmAUC{} against Mnemos-P, and the
clocked flush takes the same arm to \REthreeMLlmAUC.

\paragraph{The unmodified mem0 library}
We ran \textbf{mem0 2.2.1} with a local on-disk Qdrant collection and \texttt{\REeighteenModel}
making the add/skip decision, the adversary seeing only whether the collection's data files changed.
A fact already stored produces a vector-store write on \REeighteenMemberWrite{} of probes against
\REeighteenNonmemberWrite{} for a novel one, an AUC of \REeighteenAUC, stronger than our
mem0-style pipeline because mem0's extraction canonicalizes a repeated fact into a form its store
already holds. Since \texttt{gocryptfs} mirrors these files block for block, whether a write occurs
is the same signal with or without the encryption layer. The oracle is present in a widely used
library as shipped.

\subsection{The Dual-Epoch Game}
\label{sec:eval-game}

\begin{table*}[t]
\centering
\caption{The dual-epoch game against the compiled engines, 1000 instances per
bit, evaluated on a held-out half. \emph{Learned} is a gradient-boosted classifier over snapshot
and trace features; \emph{hand} is one named statistic per arm, oriented on the training half.
A test bounds leakage from below, so the upper end of the interval is the ceiling for these
attackers and not a certificate.}
\label{tab:e5}
\small
\setlength{\tabcolsep}{4pt}
\begin{tabular}{lrrrrl}
\toprule
Configuration & Admissible & Refusals & Learned AUC [95\% CI] & Hand AUC [95\% CI] & Verdict \\
\midrule
P shipped & yes & 0 & 0.50\,[0.47, 0.54] & 0.52\,[0.48, 0.56] & at chance \\
P no S1 (growing regions) & yes & 0 & 1.00\,[1.00, 1.00] & 1.00\,[1.00, 1.00] & distinguishable \\
P no S2 (unsealed slots) & yes & 0 & 1.00\,[1.00, 1.00] & 1.00\,[0.99, 1.00] & distinguishable \\
P no S3 (per-slot writes) & yes & 0 & 0.49\,[0.45, 0.53] & 0.50\,[0.46, 0.53] & at chance \\
P identity permutation & yes & 0 & 0.52\,[0.49, 0.56] & 0.51\,[0.47, 0.54] & at chance \\
P inadmissible, padded refusal & \textbf{no} & 7659 & 0.47\,[0.44, 0.51] & 0.53\,[0.49, 0.56] & at chance \\
P inadmissible, silent refusal & \textbf{no} & 7659 & 1.00\,[1.00, 1.00] & 0.99\,[0.99, 0.99] & distinguishable \\
M shipped & yes & 0 & 0.50\,[0.47, 0.54] & 0.51\,[0.47, 0.55] & at chance \\
M content placement & yes & 0 & 0.49\,[0.45, 0.52] & 0.49\,[0.45, 0.52] & at chance \\
P first-free slot (S3 kept) & yes & 0 & 0.51\,[0.48, 0.55] & 0.50\,[0.50, 0.50] & at chance \\
P no S3, first-free slot & yes & 0 & 1.00\,[1.00, 1.00] & 1.00\,[1.00, 1.00] & distinguishable \\
\bottomrule
\end{tabular}
\end{table*}


\Cref{tab:e5} plays the game of \cref{def:dualepoch} against the compiled engines, with
admissibility \emph{verified} (every admissible arm asserts a zero refusal counter), held-out
evaluation, and a hand-built statistic beside the learned classifier in every arm. Removing the
keyed permutation, which \cref{thm:main} does not require, leaves the game at chance once histories
are genuinely admissible; what made an identity permutation look distinguishable was the refusal
channel. Removing S1 or S2 is perfectly distinguishable. The inadmissible arms show the refusal
channel and its closure by padded refusal.

Three arms test \cref{lem:slot}, registered before they ran. Writing only the placed slot, drawn
uniformly, is at chance, \REfiveSlotRandom; so is whole-region rewriting with lowest-free-slot
placement, \REfiveRegionFirstFree. Removing both, so the index of the one slot written is the
region's load, gives \REfiveSlotFirstFree. Either protection suffices, so S3 is a condition on the
pair (placement rule, write granularity).

\subsection{The Reactive Game}
\label{sec:eval-reactive}

In the game of \cref{def:dualepoch} Mnemos-M and a variant of it that places by content are both at
chance, because the adversary supplies every operation. The reactive game separates them. We run it
with a deduplicating agent that writes a supplied fact only if no near-duplicate is stored, and
report the held-out AUC of the strongest trace statistic per design. The registered prediction was
that only the clock-only arms sit at chance.

The shipped memory-resident store is at chance, \REtwentyMSweep, and so is its whole-region-flush
variant, \REtwentyMRegions. Flushing when the agent writes is fully distinguished,
\REtwentyMOnDemand: the occurrence of a flush is the deduplication bit. Placing by content under a
clock is distinguished too, \REtwentyMContent, because a reactive write lands in the retrieved
fact's region. Mnemos-P, dual-epoch secure by \cref{thm:main}, is fully distinguished,
\REtwentyPShipped. This is the empirical form of $\mathrm{DE}\not\Rightarrow\mathrm{RDE}$, and the
reason the positive result for Mnemos-M is \cref{thm:reactive} rather than a row of \cref{tab:e5}.

\subsection{The Byte-Level Auditor}
\label{sec:eval-audit}

An auditor diffs the store byte by byte around every operation of both sealed engines and requires
each changed range to be named by an emitted event and each event to match a real change. For
Mnemos-P it confirms that a flush rewrites exactly $w+1$ whole regions, that searches, deletes and
sub-batch puts write nothing until the next flush, that a refused put still rewrites its region, and
that every byte of a rewritten region changes. For Mnemos-M it confirms that a search writes and
reads nothing, and that a sweep flush writes exactly the cursor's slot run and the journal, the
same byte ranges whether or not a record was carried. Each check fails under the ablation that removes
its mechanism. The auditor also found a defect no functional test can: rewrites owed to refused
writes were queued behind accepted ones, so a refusal-heavy history built a backlog whose length
distinguished it.

\subsection{Channels That Reach the Earlier Cohort}
\label{sec:eval-cohorth}
\label{sec:eval-write}
\label{sec:eval-multisnapshot}
\label{sec:eval-refusal}

\Cref{thm:main} closes the snapshot channel. Three channels reach cohort H anyway.

\begin{figure*}[t]
\centering
\includegraphics[width=\linewidth]{figures/fig_budget.pdf}
\caption{Exposure of the pre-compromise cohort against the adversary's induction budget, on a
timestamped stream compromised at a month boundary, for the write observer (left) and the
read-write observer (right); bands are 95\% intervals over twenty trials. The growing layout's
snapshot discloses which region holds each earlier record. For the sealed designs the harness
supplies that membership, so their curves are what the labels would be worth to an adversary that
obtained it elsewhere, for instance from version diffs. Mnemos-RS is measured for the write
observer, since its read-side channel is result-dependent reads. The supplementary material
tabulates every point.}
\label{fig:budget}
\end{figure*}

\paragraph{The write channel}
Frequency analysis labels regions with the topics of \emph{post-compromise} writes. Against a
growing layout the labels transfer to the earlier cohort; against a sealed layout they apply to a
cohort-H record only if another channel supplies its region (\cref{fig:budget}). The two adversary
classes separate here too: against Mnemos-P a read observer recovers the keyed permutation from
probe co-occurrence (\cref{sec:eval-permutation}) and labels partitions from public text with no
induced writes, and a write observer does neither. Against Mnemos-M neither learns anything.

\paragraph{Multiple versions}
Version diffs supply that membership. The region holding a pre-$\tau$ record is inside the diff at
every granularity tested; granularity sets how precisely it is named, down to the $1/B$ null when the
diff covers the whole file. Re-encrypting the whole file at a session boundary makes a
cross-session diff total and drives the precision to exactly the null.

\paragraph{The composed attack}
We ran the full chain, version diffs to membership, induced writes to labels, composed to a cohort-H
topic, on Mnemos-P at $g \in \{1, 2, 4, 16\}$. The registered prediction, that the composition
would reach cohort H, \emph{failed}. On agent memories the induced-label map does not localize
topics on the sealed, padded layout: even an oracle handed each cohort-H record's exact region scores
\REseventeenOracle, consistent with the Mnemos-P rows of \cref{tab:e2}, and the composed attack,
choosing among the \REseventeenCandOne{} candidate regions a $g=1$ diff names, scores
\REseventeenComposedOne. The multi-version channel is real, but on this corpus it does not compose
into topic recovery, because the label map it would feed is weak on the partitioned design.

\paragraph{Refusals}
An adversary that induces writes into a region and watches whether it is rewritten measures
\cref{prop:capacity} directly. A refusal that writes nothing rewrites a full region on
\REnineSilentFull{} of induced writes against \REnineSilentOpen{} for an open one, an AUC of
\REnineSilentAUC; padded refusal closes it, \REninePaddedAUC. Under a skewed corpus the partitioned
design refuses its first write at \REnineUtilP\% of nominal capacity, while Mnemos-M, with only a
global bound, reaches \REnineUtilM\%. Sweeping $B \in \{64,128,256\}$, $C \in \{16,32,64\}$ and the
Zipf skew of the topic workload (registered experiment E23), Mnemos-P's first refusal comes at
\REtwentythreePWorst\% to \REtwentythreePBest\% of nominal capacity, earlier as $B$ grows and as the
skew rises (\REtwentythreePSkewLow\% to \REtwentythreePSkewHigh\% from Zipf 0.8 to 1.2), because the
first partition to fill is the most frequent topic's; Mnemos-M's never comes before
\REtwentythreeMWorst\%. A partitioned store that must keep every write is provisioned for its most
frequent topic; padded refusal trades that provisioning for lost writes, a cost in availability,
not confidentiality.

\subsection{Read-Side Channels, for the RW Class Only}
\label{sec:eval-rw}
\label{sec:eval-search}
\label{sec:eval-query}
\label{sec:eval-permutation}

These results apply to the read-write observer, not to the write observer of the anchoring
instantiations, which sees no reads.

\paragraph{Search and deletion}
The record-store design issues \RSearchReadsRS{} result-dependent reads over a search trace,
attaching a query topic to \RSearchPositionsRS{} record positions, and names the region of a deleted
cohort-H record immediately in \RDeleteNamedRS{} of \RDeleteN{} deletions. The partitioned design
issues \RSearchReadsP{} such reads and names \RDeleteNamedP{} regions; Mnemos-M's search does not
touch the disk.

\paragraph{Query equality}
Over \REfifteenRepeats{} independent stores, with the adversary's overlap threshold fixed on half the
queries and scored on the rest, the shipped randomized superset defeats a test for exact equality of
read sets (Youden's $J = \REfifteenExactStwo$) but not a threshold on their overlap,
$J = \REfifteenThreshStwo$, because repeated searches share every genuine probe. Grouped probing
leaks pairwise equality at group granularity, $J = \REfifteenThreshGrouped$. Sweeping each knob
(registered experiment E22, 20 stores per point) traces the frontier. The spread leaks
$J = \REtwentytwoSpreadOne$ at $s = 1$, $\REtwentytwoSpreadFour$ at $s = 4$ and
$\REtwentytwoSpreadEight$ at $s = 8$, eight times the genuine reads; grouping leaks
$\REtwentytwoGroupTwo$ at group size 2 and $\REtwentytwoGroupThirtyTwo$ at 32
(\REtwentytwoGroupAmpThirtyTwo$\times$ the reads). Neither dominates: at twice the reads grouping
leaks less ($\REtwentytwoGroupEight$ against $\REtwentytwoSpreadTwo$), from four times upward the
spread does, and no point reaches chance. Closing the channel requires oblivious access.

\paragraph{Permutation recovery}
Genuine probes are the nearest partitions to a query, so regions read together are geometric
neighbours, and matching the observed co-occurrence graph to the one the public table predicts
recovers the permutation from ordinary search traffic, with no induced writes. Without probe
randomization, an approximate quadratic-assignment solver refined by pairwise swaps places
\REsevenExactCISone{} of the $B$ regions exactly over \REsevenRepeats{} independent stores (the
solver alone reaches \REsevenFaqSone, so the number is a lower bound). The randomized superset
degrades this to \REsevenExactStwo, and grouped probing defeats it, \REsevenExactGrouped{} against a
$1/B$ null of 0.016. The permutation protects against a write observer holding a public table, not
against a read observer.

\subsection{Cost}
\label{sec:eval-perf}

We measure the three designs against the in-tree Rust plaintext IVF, which shares
\texttt{QuantVec} and the partitioner with the encrypted engines, at five sizes, with writes timed
in the sparse regime where an agent writing seconds apart pays a whole flush per record (the
supplementary material plots every measure against $n$). Exhaustive in-memory ranking is
\REtenSpeedupSmall{} times faster than the partitioned design at \REtenNSmall{} records and
\REtenSpeedupLarge{} times as fast at \REtenNLarge, crossing it between \REtenCrossLo{} and
\REtenCrossHi{} records, because it is linear in $n$. Probing logical partitions in memory removes the inversion without reopening the write
channel: at eight partitions scored per query, Mnemos-M with its in-memory index is
\REtenIvfSpeedupMin{} to \REtenIvfSpeedupMax{} times faster than Mnemos-P, \REtenIvfSearchLarge~ms
against \REtenPSearchLarge~ms at the largest size, because it reads and decrypts nothing. It pays at
unlock, \REtenUnlockIvfLarge~s against \REtenUnlockMLarge~s at \REtenNLarge{} records, and in
resident memory, \REtenResidentLarge~MB. A sparse write costs \REtenPutMLarge~ms against
\REtenPutPLarge~ms. The supplementary material tabulates every size and prices fetching results
once reads are visible (by position, by full scan, and Path ORAM as a stated model).

\paragraph{What the clocked flush writes}
A clocked flush is paid whether or not the agent wrote. Rewriting $w+1$ whole regions each tick is
\REtwentyoneRegionMiB~MiB a flush at the shipped geometry, \REtwentyoneRegionGBDay~GB a day at a
$500$~ms clock. The sweep flush, in the manner of write-only
ORAM~\cite{roche2017detwooram,aviv2017oblivisync,blass2014hive}, rewrites the next fixed run of slots
after a cursor plus a fixed-size sealed journal that holds staged records until the cursor reaches a
free slot. It writes \REtwentyoneSweepKiB~KiB per flush whether the flush carries a record or not,
and at $\Delta = 10$~s, \REtwentyoneSweepGBDay~GB a day: roughly \REtwentyoneRatio${\times}$ less, and
about \REtwentyoneDriveYears{} years of continuous operation against a $600$~TBW rating. Locking
Mnemos-M rewrites nothing extra; locking Mnemos-P rewrites its whole file,
\REtwentyonePLockMiB~MiB, which a synchronization client uploads in full each session.

\paragraph{Above the memory limit}
Against the write observer W that anchors the threat model, the partitioned design is not the
right tool above the memory limit. Mnemos-M/disk (\cref{sec:mnemos-m}) keeps only the position map
and logical index in memory and leaves headers and payloads on disk, written only by the sweep, so
its writes remain a function of the clock and every write-side channel stays closed against W at
any size. A read observer RW sees its search access pattern, which only oblivious search closes;
\cref{sec:background} compares that cost against recent disk-oblivious ANN and private-memory
systems~\cite{rathee2026onyx,kaviani2026opal,zhu2025compass}. The partitioned design remains the
construction \cref{thm:main} is proved for and the point of comparison, not the recommended
deployment.

\subsection{Retrieval Quality}
\label{sec:eval-recall}

Retrieval is measured against BEIR relevance judgments on the \emph{full} corpora (truncating FiQA
to 5{,}000 documents, under a tenth of its \REelevenFiqaSize, inflates plaintext nDCG@10 well above
\REelevenFiqa), against a plaintext IVF at the same probe width, and counting capacity refusals
during indexing (\REelevenRefusals). Quantisation to int8 costs at most \REelevenWorstQuant{}
nDCG@10 on any domain, and fixed probing up to \REelevenProbeCostMax. Mnemos-M's in-memory index,
scoring the same eight partitions as Mnemos-P at $s = 2$, gains \REelevenIvfGainMin{} to
\REelevenIvfGainMax{} on every domain, since all eight are the query's nearest where half of
Mnemos-P's are cover.

Grouped probing, the one mechanism that resists permutation recovery against a read observer, is
also the one that costs retrieval: reading a whole disjoint group per query drops nDCG@10 from
\REelevenGroupedSScifact{} to \REelevenGroupedScifact{} on SciFact and from \REelevenGroupedSFiqa{}
to \REelevenGroupedFiqa{} on FiQA against the $s=2$ spread. That price is why the memory-resident
design, which removes the read channel by keeping search off the disk, is the one we recommend.

\subsection{Operating-System Artifacts}
\label{sec:eval-artifacts}

The audit of \cref{sec:artifacts} is run, not asserted. On the evaluation machine it recovers
change-journal records naming the store's files with timestamps, enumerates shadow copies and
synchronized trees, reports the hibernation file, and verifies that after a lock both sealed engines
refuse to search. We grade each artifact on a five-rung ladder: it exists; it contains data relevant
to the store; the adversary can obtain it; it is exploitable for the historical attack; it suffices
alone. The change journal reaches rung~2: it records when the files changed, not a region diff.
Shadow copies and file-history snapshots reach rung~4, as full prior versions that give the
multi-version channel at the snapshot schedule's granularity. Retained flash pages reach rung~4 for
Mnemos-P, whose whole-region rewrites land at fixed offsets, and rung~1 for Mnemos-M. The
hibernation image is the one rung-5 artifact, against key custody rather than the layout, and is
handled as a deployment requirement. The study is Windows-specific (USN journal, VSS, CNG); the
Linux counterparts, filesystem journals, LVM or Btrfs snapshots and swap, occupy the same rungs but
are not measured.

\subsection{Where the Information Goes}
\label{sec:eval-ledger}

\begin{table*}[t]
\centering
\caption{Every channel in one unit. \emph{Bits} is a Fano lower bound on the mutual information between the secret and what the adversary observes, per observation of the unit the channel is about (a record's topic, a region's partition, one game instance), taken at the lower end of the experiment's 95\% interval so that it holds at the experiment's confidence; $^\ast$ marks a result with no interval, for which only the point bound exists. A test bounds leakage from below, so for a \emph{closed} channel the informative number is the last column: the smallest effect, in the experiment's own unit, it had 80\% power to detect. The supplement gives every row.}
\label{tab:ledger}
\footnotesize
\setlength{\tabcolsep}{4pt}
\begin{tabular}{llllrl}
\toprule
Channel & Design & Class & Status & Bits $\geq$ & Detectable $\geq$ \\
\midrule
structure transfer, historically aware adversary & growing layout & W & \textbf{open} & 0.815 & -- \\
snapshot + post-$\tau$ writes (the game) & Mnemos-P & W & closed & 0.000 & 0.052 (AUC) \\
snapshot + post-$\tau$ writes (the game) & Mnemos-M & W & closed & 0.000 & 0.051 (AUC) \\
refusal, silent & Mnemos-P ablation & W & \textbf{open} & 0.682 & -- \\
refusal, silent & Mnemos-P & W & \textbf{open} & 0.763 & -- \\
refusal, padded & Mnemos-P & W & closed & 0.000 & 0.051 (AUC) \\
refusal, padded & Mnemos-P & W & closed & 0.000 & undetermined \\
multi-version diff, g=1 & Mnemos-P & W & \textbf{open} & 0.524 & -- \\
multi-version diff, across a session, rewrite\_at\_lock on & Mnemos-P & W & closed & 0.000 & undetermined \\
permutation recovery (spread\_s1) & Mnemos-P & RW & \textbf{open} & 4.905 & -- \\
permutation recovery (spread\_s2) & Mnemos-P & RW & \textbf{open} & 1.827 & -- \\
permutation recovery (grouped) & Mnemos-P & RW & \textbf{open} & 0.000 & -- \\
query equality (spread\_s2), overlap threshold & Mnemos-P & RW & \textbf{open} & 0.890 & -- \\
query equality (grouped), overlap threshold & Mnemos-P & RW & \textbf{open} & 0.566 & -- \\
dedup membership (threshold) & Mnemos-P & W & \textbf{open} & 0.189 & -- \\
dedup membership (threshold) & Mnemos-M & W & closed & 0.000 & undetermined \\
dedup membership (llm) & Mnemos-P & W & \textbf{open} & 0.003 & -- \\
dedup membership (llm) & Mnemos-M & W & closed & 0.000 & undetermined \\
\bottomrule
\end{tabular}
\end{table*}


\Cref{tab:ledger} puts the channels on one axis. Where an experiment recorded a held-out accuracy,
Fano's inequality with the secret's empirical entropy turns it into a lower bound on the mutual
information between the secret and the adversary's view; an AUC is reported as an AUC and is not
converted. At a spread of one, permutation recovery carries at least \RLedgerPermBits{} of the
$\log_2 B = \RLedgerPermBitsMax$ bits of each region's partition, and grouped probing takes it to
zero; the multi-version channel depends on capture granularity and vanishes across a session once
the file is re-encrypted at lock. For a row where nothing was detected, the informative number is
the last column, the smallest effect the experiment had 80\% power to detect: ``not detected'' is a
bounded statement about this attacker at this sample size, not a certificate.
